\chapter{Centralized Configuration: Config Server}
\label{ch:configserver}

\section{The problem}

Seven services each need database URLs, Kafka addresses, JWT settings,
mail hosts. Copy that into seven \code{application.yml} files and every
change becomes seven edits, seven rebuilds, and the guarantee that two
copies will eventually disagree. The classic Spring Cloud answer:
one small HTTP service whose only job is handing each application its
configuration at startup.

\section{How Spring Cloud Config works}

The server exposes configuration over a plain HTTP contract:

\begin{lstlisting}
GET /{application}/{profile}
GET /course-service/default      -> the config for course-service
\end{lstlisting}

The client side is the \code{spring.config.import:} \code{configserver:...}
line you saw in Chapter~1. During startup --- \emph{before} most beans are
created --- the client fetches its document and merges it into the
Environment as a property source that outranks the packaged YAML but
loses to environment variables. The resolution order of
Section~\ref{sec:property-order} is doing all the work; the config-server
is just one more source in the stack.

\section{The project's server, annotated}

\begin{lstlisting}[language=yaml]
server:
  port: 8888                       (*@\co{1}@*)
spring:
  application:
    name: config-server
  profiles:
    active: native                 (*@\co{2}@*)
  cloud:
    config:
      server:
        native:
          search-locations: classpath:/configurations   (*@\co{3}@*)
management:
  endpoints:
    web:
      exposure:
        include: health            (*@\co{4}@*)
\end{lstlisting}

\begin{description}
  \item[\co{1}] 8888 is the conventional config-server port, the way 8761
  is Eureka's. Convention saves documentation.
  \item[\co{2}] The \code{native} profile means ``serve files from a
  filesystem/classpath location'' instead of the default git backend.
  \item[\co{3}] The served files are \emph{inside the config-server's own
  jar}: \code{configurations/course-service.yml},
  \code{auth-service.yml}, and so on --- selected by matching file name to
  the requesting service's \code{spring.application.name}.
  \item[\co{4}] Only \code{health} is exposed --- and this is the endpoint
  the Compose healthcheck and the Kubernetes probes rely on.
\end{description}

\begin{decision}[native classpath vs.\ git backend]
The enterprise-standard backend is a git repository: config changes are
commits --- audited, reviewed, revertible without rebuilding anything.
The cost is a second repository and credentials management. For a
single-developer project whose config changes ship with code changes
anyway, classpath-native keeps everything in one repo and one pipeline.
The trade traded away: changing a value now requires rebuilding the
config-server image. Know which side of that trade your employer sits on.
\end{decision}

\begin{pitfall}[a service starts with defaults and fails oddly]
Symptom: auth-service boots but every request fails with connection
errors to \code{localhost:5432}. Cause: the config-server was unreachable
at startup, and \code{optional:configserver:} let the service boot from
its packaged defaults --- which point at localhost. Because the failure
appears minutes later and far from the cause, it reads as a database
problem. Check: the startup log prints whether remote config was fetched.
The Compose file prevents this ordering problem with
\code{depends\_on: condition: service\_healthy}; Kubernetes has no
startup ordering (Chapter~15), so services must tolerate a late
config-server by crashing and restarting until it appears.
\end{pitfall}

\section{Startup order and the bootstrap dance}

The config-server creates a hidden sequencing constraint across the whole
system: every other service wants it up first. You will meet the
consequences three times in this book:

\begin{itemize}
  \item Compose encodes the order explicitly with \code{depends\_on} and a
  \code{curl}-based healthcheck (Chapter~12).
  \item Kubernetes refuses to encode startup order at all; the system
  converges by retry (Chapter~15).
  \item The Jenkins pipeline deploys all manifests in one
  \code{kubectl apply -k} and lets rollouts settle concurrently
  (Chapter~23).
\end{itemize}

\section{Refreshing config without restart}

Out of the box, a client reads its config once at startup. Spring Cloud
offers \code{/actuator/refresh} (re-fetch on demand, per instance,
applies to \code{@RefreshScope} beans) and Spring Cloud Bus (broadcast
refresh over a message broker). This project uses neither: on Kubernetes
a config change ships as a new image or env change, and the Deployment's
rolling update (Chapter~15) restarts pods anyway. That is the modern
answer to a question the config-server was originally designed around.

\section{Outside the project}

A typical git-backed enterprise setup, for contrast:

\begin{lstlisting}[language=yaml]
spring:
  cloud:
    config:
      server:
        git:
          uri: https://git.company.com/platform/config-repo
          default-label: main
          search-paths: '{application}'
\end{lstlisting}

Config lives per-service in folders of a dedicated repo; the label is a
branch, so \code{default-label: release-42} pins an environment to a
release. The second widespread pattern skips the config-server entirely:
Kubernetes ConfigMaps and Secrets mounted as env vars --- which is
exactly where this project is heading (Chapters~16 and 26); the
config-server survives here mostly as a working example of the classic
stack.

\begin{quiz}
\begin{enumerate}
  \item In which order do these sources win for
  \code{spring.datasource.url}: packaged YAML, config-server document,
  Deployment env var?
  \item What does \code{profiles.active: native} change about where the
  server finds its documents, and what does the git backend buy that
  native cannot?
  \item Why is \code{optional:} in \code{spring.config.import} a
  double-edged sword? Describe the failure mode it permits.
  \item The config-server's own config cannot come from a config-server.
  Where must it live, and what general bootstrapping principle is that?
\end{enumerate}
\end{quiz}
